\chapter*{Lời mở đầu}

Trong thập niên vừa qua, các ứng dụng nhắn tin đã trở thành phương tiện giao tiếp chủ yếu của hàng tỷ người dùng trên thế giới. Cùng với sự phát triển mạnh mẽ đó, mức độ nhạy cảm của dữ liệu được trao đổi qua các ứng dụng này cũng ngày một lớn, từ thông tin cá nhân, giao dịch tài chính cho đến các tài liệu nội bộ của tổ chức. Nhu cầu bảo mật thông tin trong quá trình truyền tải trở thành một yêu cầu bắt buộc, và mã hóa đầu cuối đã dần trở thành tiêu chuẩn trong thiết kế các nền tảng nhắn tin hiện đại như Signal, WhatsApp hay iMessage.

Tuy nhiên, phần lớn các hệ mã hóa khóa công khai đang được sử dụng hiện nay đều dựa trên các bài toán khó truyền thống như phân tích số nguyên lớn trong RSA hay bài toán logarit rời rạc trên đường cong elliptic trong ECDH. Sự xuất hiện của máy tính lượng tử cùng thuật toán Shor đã đặt ra nguy cơ rằng toàn bộ các giao thức trao đổi khóa dựa trên nền tảng này sẽ không còn an toàn trong tương lai không xa. Mối đe dọa không chỉ giới hạn ở thời điểm máy tính lượng tử thực sự xuất hiện, mà còn hiện hữu ngay lúc này qua mô hình tấn công “harvest now, decrypt later”, trong đó kẻ tấn công thu thập trước các bản mã và chờ đến khi công nghệ đủ mạnh để giải mã.

Nhằm giải quyết thách thức đó, Viện Tiêu chuẩn và Công nghệ Quốc gia Hoa Kỳ (NIST) đã tổ chức chương trình chuẩn hóa mật mã hậu lượng tử kéo dài nhiều năm, và đến năm 2024 đã công bố các chuẩn chính thức bao gồm ML-KEM, được phát triển dựa trên thuật toán Kyber. Đây là thuật toán trao đổi khóa dạng Key Encapsulation Mechanism, xây dựng trên bài toán Module Learning With Errors trong mật mã trên lưới (lattice), được thiết kế để kháng lại cả tấn công cổ điển lẫn tấn công bằng máy tính lượng tử.

Xuất phát từ bối cảnh thực tế cùng mong muốn tìm hiểu sâu hơn về mật mã học hiện đại và các công nghệ bảo mật mạng, em đã chọn đề tài “Xây dựng giao thức nhắn tin an toàn tích hợp giải pháp mật mã hậu lượng tử” để làm đồ án tốt nghiệp. Đề tài hướng tới việc khảo sát nguyên lý của các thuật toán mật mã hậu lượng tử, thiết kế một kiến trúc nhắn tin đầu cuối có bắt tay khóa hybrid hậu lượng tử, và hiện thực hóa kiến trúc này thành một ứng dụng desktop có các chức năng chính của một nền tảng nhắn tin thực tế.

Nội dung chính của đồ án bao gồm 3 chương ngoài phần mở đầu và kết luận, cụ thể:

\begin{itemize}
  \item \textbf{Chương 1: Kiến thức cơ sở}: Chương này trình bày các kiến thức nền tảng về mã hóa đầu cuối trong ứng dụng nhắn tin, mật mã đối xứng hiện đại, mối đe dọa của máy tính lượng tử đối với mật mã khóa công khai truyền thống, nguyên lý của mật mã hậu lượng tử và thuật toán ML-KEM, cùng các công nghệ được sử dụng để triển khai ứng dụng.
  \item \textbf{Chương 2: Phân tích và thiết kế hệ thống}: Chương này tập trung khảo sát các ứng dụng nhắn tin bảo mật hiện hành, phân tích yêu cầu, đề xuất kiến trúc tổng thể và trình bày thiết kế các thành phần cốt lõi của hệ thống, bao gồm cơ chế E2EE, bắt tay khóa kiểu PQXDH, Double Ratchet, group theo epoch, identity transparency, cơ sở dữ liệu, client và server.
  \item \textbf{Chương 3: Cài đặt và đánh giá hệ thống}: Chương này trình bày cách triển khai hệ thống bằng Docker, mô tả các luồng sử dụng chính, quan sát dữ liệu lưu trên server, đánh giá bảo mật, kiến trúc phần mềm, khả năng sử dụng, hiệu năng, giới hạn còn lại và các hướng phát triển tiếp theo.
\end{itemize}

Cuối cùng, em xin chân thành cảm ơn các thầy cô trong khoa Toán-Tin đã tận tình giảng dạy và trang bị cho em những kiến thức quý báu trong suốt quá trình học tập. Đặc biệt, em xin gửi lời cảm ơn sâu sắc đến thầy Trần Anh Tú đã luôn quan tâm, hướng dẫn và tạo điều kiện giúp em hoàn thành đồ án này.
